\chapter{Metodología}

Para el desarrollo del estudio simplificado de conexión, se realizaron los estudios de flujo de carga y cortocircuito en el NODO 3272966. A continuación se describen los análisis realizados.

\section{Flujo de Carga}

El análisis de flujo de carga tiene como objetivo determinar el desempeño del sistema eléctrico ante la conexión de la nueva generación. Específicamente, busca evaluar si en condiciones de estado estable, la incorporación del proyecto fotovoltaico genera dificultades operativas en la red, tales como:

\begin{itemize}
    \item Sobrecargas en líneas y transformadores
    \item Violaciones de los límites de tensión establecidos
\end{itemize}

Este análisis busca garantizar que el suministro de energía se mantenga acorde con las exigencias de calidad, seguridad y balance del sistema, conforme a lo dispuesto en:

\begin{itemize}
    \item Resolución CREG 025 de 1995 (Código de Redes), con sus modificaciones vigentes
    \item Resolución CREG 174 de 2021 (Lineamientos para conexión de autogeneradores a pequeña escala)
\end{itemize}

\subsection{Límites Operativos Considerados}

\begin{enumerate}
    \item \textbf{Tensiones de estado estable:} Se consideran aceptables cuando se encuentran entre el 90\% y el 110\% de la tensión nominal del barraje correspondiente.
    
    \item \textbf{Cargabilidad de líneas y transformadores:} Se evaluó con base en la corriente nominal, asumiendo que estos elementos pueden operar hasta el 100\% de su capacidad nominal sin constituir una sobrecarga.
\end{enumerate}

\subsection{Metodología de Análisis}

A través de este análisis se determinan:

\begin{itemize}
    \item Las tensiones en las barras del sistema
    \item La distribución de los flujos de potencia activa y reactiva
    \item Las pérdidas técnicas asociadas a la operación de la red
\end{itemize}

Se comparan dos escenarios:
\begin{enumerate}
    \item Sin generación fotovoltaica
    \item Con incorporación del generador fotovoltaico
\end{enumerate}

El análisis se realiza para las siguientes franjas horarias:
\begin{itemize}
    \item 9:00 a.m.
    \item 12:00 p.m.
    \item 3:00 p.m.
\end{itemize}

\section{Análisis de Cortocircuito}

El análisis de cortocircuito se realizó con base en la información de red suministrada por el operador. Se calcularon las corrientes de cortocircuito trifásico (L-L-L) y monofásico (L-G) para los nodos aledaños al proyecto de generación fotovoltaica.

\subsection{Normas Aplicadas}

Los cálculos se realizaron conforme a la norma IEC 60909, según la cual las tensiones de prefalla en las barras corresponden al 110\% de la tensión nominal.

\subsection{Contribuciones de Cortocircuito}

\begin{enumerate}
    \item \textbf{Equivalente de red:} Se consideró el aporte de cortocircuito definido por el equivalente de red en la subestación CONUCOS, a 13,8 kV.
    
    \item \textbf{Microinversores:} Se asumió que la corriente de cortocircuito simétrica equivale a 1,5 veces la corriente nominal del microinversor, conforme a criterios habitualmente empleados para este tipo de tecnología.
\end{enumerate}

\subsection{Parámetros Calculados}

Como resultado de este análisis se obtienen:

\begin{itemize}
    \item \textbf{Corrientes de interrupción simétrica de cortocircuito ($I_k$):} Valor r.m.s. de la componente simétrica de la corriente de cortocircuito en el instante de separación de los contactos del dispositivo de interrupción.
    \item \textbf{Corrientes de cresta ($I_p$):} Valor máximo instantáneo esperado de la corriente de cortocircuito.
\end{itemize}

\section{Modelado del Sistema}

\subsection{Datos de Entrada}

Para modelar la red equivalente propuesta por el operador de red, se utilizaron los datos suministrados en el archivo ``Datos Circuito 10502.xlsx'' con el cual se modeló todo el circuito.

\subsection{Características del Circuito}

La red equivalente del sistema de distribución se modeló a partir de:

\begin{itemize}
    \item Parámetros eléctricos de las líneas
    \item Topología del circuito 10502
    \item Registro horario de la demanda en el alimentador
\end{itemize}

\subsection{Ajustes de Demanda}

Con base en la información disponible, se caracterizó el comportamiento de la corriente en la cabecera del circuito a lo largo del día, y se ajustó la demanda del sistema para las tres franjas horarias evaluadas:

\begin{itemize}
    \item 9:00 a.m.
    \item 12:00 p.m.
    \item 3:00 p.m.
\end{itemize}

\subsection{Elementos Modelados}

\subsubsection{Transformador de Conexión}

El transformador de conexión se modeló con los parámetros presentados en la tabla correspondiente:

\begin{itemize}
    \item Potencia: 150 kVA
    \item Relación: 13,2 kV / 220 V
    \item Grupo de conexión: Dyn5
    \item Impedancia de cortocircuito: 4,9\%
\end{itemize}

\subsubsection{Líneas y Conductores}

Las características técnicas de los conductores empleados en el modelado corresponden al tramo que conecta el nodo 3272869 con el nodo 3272966:

\begin{itemize}
    \item Cable tipo CU 2 XLPE
    \item Longitud: 3 metros
    \item Corriente nominal: 155 A
\end{itemize}

\section{Criterios de Aceptación}

Los resultados del análisis se comparan contra los siguientes criterios:

\begin{table}[H]
    \centering
    \caption{Criterios de aceptación para el análisis}
    \begin{tabular}{ll}
    \toprule
    \textbf{Parámetro} & \textbf{Criterio} \\
    \midrule
    Tensión en barras & 0.90 - 1.10 p.u. \\
    Cargabilidad de líneas & $\leq$ 100\% \\
    Cargabilidad de transformadores & $\leq$ 100\% \\
    Aumento de cortocircuito & $<$ 1\% \\
    \bottomrule
    \end{tabular}
\end{table}
